% Mức: NB
% Mức: NB
% ID: L12C1B3-121-TN
% ID: L12C1B3-121-TN
% Mức: NB
% ID: L12C1B3-121-TN
% ID: L12C1B3-121-TN
% Mức: NB
%=== Câu 1 ===%
\dangbt{Khái niệm nhiệt độ}
\begin{ex}
% ID: L12C1B3-121-TN
% Mức: NB
	Các vật \textbf{không} thể có nhiệt độ thấp hơn
	\choice
	{$2,0\ K$}
	{$0^\circ C$}
	{$100\ K$}
	{\True $-273,15^\circ C$}
	\loigiai{
	Nhiệt độ của các vật không thể nhỏ hơn $0\ K = -273,15^\circ C$.
	}
\end{ex}

% Mức: NB
% Mức: NB
% ID: L12C1B3-122-TN
% ID: L12C1B3-122-TN
% Mức: NB
% ID: L12C1B3-122-TN
% ID: L12C1B3-122-TN
% Mức: NB
%=== Câu 2 ===%
\begin{ex}
% ID: L12C1B3-122-TN
% Mức: NB
	Phát biểu nào sau đây đúng?
	\choice 
	{Khi một vật tỏa nhiệt ra môi trường thì nội năng của vật tăng lên}
	{Độ biến thiên nội năng của một vật là độ biến thiên nhiệt độ của vật đó}
	{Nội năng là phần năng lượng vật nhận được hay mất đi trong quá trình truyền nhiệt}
	{\True Nội năng của vật phụ thuộc vào nhiệt độ và thể tích của vật}
	\loigiai{
		\begin{itemize}
			\item Phát biểu  Nội năng của vật phụ thuộc vào nhiệt độ và thể tích của vật đúng vì nội năng của một hệ là tổng động năng và thế năng tương tác của các phân tử cấu tạo nên hệ. Động năng phân tử phụ thuộc vào nhiệt độ, còn thế năng tương tác phân tử phụ thuộc vào khoảng cách giữa các phân tử (tức là phụ thuộc vào thể tích của vật). Do đó, nội năng của vật phụ thuộc vào nhiệt độ và thể tích,,.
			\item Các phát biểu khác sai vì: Khi một vật toả nhiệt ra môi trường thì nội năng của vật giảm xuống. Phần năng lượng nhiệt mà vật nhận được hay mất đi trong quá trình truyền nhiệt được gọi là nhiệt lượng.
		\end{itemize}
	}
\end{ex}

% Mức: VD
% Mức: VD
% ID: L12C1B3-129-TN
% ID: L12C1B3-129-TN
% Mức: VD
% ID: L12C1B3-129-TN
% ID: L12C1B3-129-TN
% Mức: VD
%=== Câu 9 ===%
\dangbt{Thang nhiệt độ và nhiệt kế}
\begin{ex}
% ID: L12C1B3-129-TN
% Mức: VD
	Chiều cao của cột thủy ngân trong nhiệt kế thủy ngân thay đổi theo nhiệt độ. Ứng với hai vạch có nhiệt độ là $0^\circ C$ và $100^\circ C$ thì chiều cao của cột thủy ngân trong nhiệt kế là $2\ cm$ và $22\ cm$. Khi sử dụng nhiệt kế này để đo nhiệt độ của cơ thể của một em bé đang bị sốt thì thấy cột thủy ngân cao $9,9 \ cm$. Theo thang nhiệt Kelvin, nhiệt độ của em bé lúc này là bao nhiêu?
	\choice 
	{$321,5\ K$}
	{$305,5\ K$}
	{$327,0\ K$}
	{\True $312,5 K$}
	\loigiai{
		Ta có: $\dfrac{\Delta t_1 }{\Delta x_1} = \dfrac{\Delta t_2}{\Delta x_2} \Leftrightarrow \dfrac{100-0}{22-2} = \dfrac{t-0}{9,9 - 2} \Rightarrow t = 39,5^\circ C \Rightarrow T = 312,5\ K$.
	}
\end{ex}

% Mức: VD
% Mức: VD
% ID: L12C1B3-130-TN
% ID: L12C1B3-130-TN
% Mức: VD
% ID: L12C1B3-130-TN
% ID: L12C1B3-130-TN
% Mức: VD
%=== Câu 10 ===%
\begin{ex}
% ID: L12C1B3-130-TN
% Mức: VD
	Hình bên là dự báo thời tiết ở thành phố Hồ Chí Minh ngày $21/07/2024$. Theo thang nhiệt Kelvin, nhiệt độ cao nhất và thấp nhất theo bảng dự báo bên là
	\begin{center}
		\begin{tikzpicture}
			%\draw [dotted] (-2,-2) grid (14,4);
			%Hình 1
			\draw [rounded corners](-1.2,-1.5) rectangle (1.2,2);
			\node [star, star points=12, star point ratio=0.4, draw, fill=yellow, minimum size=1cm, scale=0.5] at (0.4,0.4) {};
			\begin{scope}[shift={(0.5,-0.5)}]
				\draw [fill = yellow, rotate=-20] (0,0.5) -- (-0.2,-0.2) -- (-0.1,-0.2) -- (-0.2,-0.2-0.7) -- (0,-0.1)--(-0.1,-0.1)--cycle ;
			\end{scope}
			\node[cloud, draw, minimum width=2cm, minimum height=1cm, fill=blue!50] at (0,0) {};
			\foreach \x/\y in {-0.5/-0.5,-0.2/-0.9,0.1/-0.7,0.
				4/-0.8,0.7/-0.4,0.9/-0.7}{
				\draw [fill = blue, rotate=-20] (\x,\y) .. controls (\x-0.2,\y-0.2) and (\x+0.2,\y-0.2) .. (\x,\y);
			}
			\draw 
			node[above] at (0,1) {$28^\circ C$}
			node[above] at (0,1.5) {10:00};
			%Hình 2
			\begin{scope}[shift={(3,0)}]
				\draw [rounded corners](-1.2,-1.5) rectangle (1.2,2);
				\node [star, star points=12, star point ratio=0.4, draw, fill=yellow, minimum size=1cm, scale=0.5] at (0.4,0.4) {};
				\begin{scope}[shift={(0.5,-0.5)}]
					\draw [fill = yellow, rotate=-20] (0,0.5) -- (-0.2,-0.2) -- (-0.1,-0.2) -- (-0.2,-0.2-0.7) -- (0,-0.1)--(-0.1,-0.1)--cycle ;
				\end{scope}
				\node[cloud, draw, minimum width=2cm, minimum height=1cm, fill=blue!50] at (0,0) {};
				\foreach \x/\y in {-0.5/-0.5,-0.2/-0.9,0.1/-0.7,0.
					4/-0.8,0.7/-0.4,0.9/-0.7}{
					\draw [fill = blue, rotate=-20] (\x,\y) .. controls (\x-0.2,\y-0.2) and (\x+0.2,\y-0.2) .. (\x,\y);
				}
				\draw 
				node[above] at (0,1) {$32^\circ C$}
				node[above] at (0,1.5) {13:00};
			\end{scope}
			%Hình 3
			\begin{scope}[shift={(6,0)}]
				\draw [rounded corners](-1.2,-1.5) rectangle (1.2,2);
				\node [star, star points=12, star point ratio=0.4, draw, fill=yellow, minimum size=1cm, scale=0.5] at (0.4,0.4) {};
				\begin{scope}[shift={(0.5,-0.5)}]
					\draw [fill = yellow, rotate=-20] (0,0.5) -- (-0.2,-0.2) -- (-0.1,-0.2) -- (-0.2,-0.2-0.7) -- (0,-0.1)--(-0.1,-0.1)--cycle ;
				\end{scope}
				\node[cloud, draw, minimum width=2cm, minimum height=1cm, fill=blue!50] at (0,0) {};
				\foreach \x/\y in {-0.5/-0.5,-0.2/-0.9,0.1/-0.7,0.
					4/-0.8,0.7/-0.4,0.9/-0.7}{
					\draw [fill = blue, rotate=-20] (\x,\y) .. controls (\x-0.2,\y-0.2) and (\x+0.2,\y-0.2) .. (\x,\y);
				}
				\draw 
				node[above] at (0,1) {$30^\circ C$}
				node[above] at (0,1.5) {15:00};
			\end{scope}
			%Hình 4
			\begin{scope}[shift={(9,0)}]
				\draw [rounded corners](-1.2,-1.5) rectangle (1.2,2);
				\node [star, star points=12, star point ratio=0.4, draw, fill=yellow, minimum size=1cm, scale=0.5] at (0.4,0.4) {};
				\begin{scope}[shift={(0.5,-0.5)}]
					\draw [fill = yellow, rotate=-20] (0,0.5) -- (-0.2,-0.2) -- (-0.1,-0.2) -- (-0.2,-0.2-0.7) -- (0,-0.1)--(-0.1,-0.1)--cycle ;
				\end{scope}
				\node[cloud, draw, minimum width=2cm, minimum height=1cm, fill=blue!50] at (0,0) {};
				\foreach \x/\y in {-0.5/-0.5,-0.2/-0.9,0.1/-0.7,0.
					4/-0.8,0.7/-0.4,0.9/-0.7}{
					\draw [fill = blue, rotate=-20] (\x,\y) .. controls (\x-0.2,\y-0.2) and (\x+0.2,\y-0.2) .. (\x,\y);
				}
				\draw 
				node[above] at (0,1) {$27^\circ C$}
				node[above] at (0,1.5) {15:00};
			\end{scope}
			\draw
			node[above] at (4.5,2) {Tp Hồ Chí Minh, ngày 21/07/2024};
		\end{tikzpicture}
	\end{center}
	\choice 
	{$246\ K$ và $241\ K$}
	{\True $305\ K$ và $300\ K$}
	{$248\ K$ và $236\ K$}
	{$321\ K$ và $306\ K$}
	\loigiai{
		Dựa vào hình, ta thấy:
		$\begin{cases}
			t_{max} = 32^\circ C &\Rightarrow T_{max} = 305\ K\\
			t_{min} = 32^\circ C &\Rightarrow T_{min} = 300\ K
		\end{cases}$.
	}
\end{ex}

% Mức: TH
% Mức: TH
% ID: L12C1B3-131-TN
% ID: L12C1B3-131-TN
% Mức: TH
% ID: L12C1B3-131-TN
% ID: L12C1B3-131-TN
% Mức: TH
%=== Câu 11 ===%
\begin{ex}
% ID: L12C1B3-131-TN
% Mức: TH
	Một nhiệt kế có phạm vi đo từ $263 \ K$ đến $1273 \ K$, dùng để đo nhiệt độ của các lò nung. Phạm vi đo của nhiệt kế này trong thang nhiệt độ Celsius là
	\choice 
	{0$^\circ C$ đến 273$^\circ C$}
	{$-20^\circ C$ đến $1200^\circ C$}
	{\True $-10^\circ C$ đến $1000^\circ C$}
	{$-12^\circ C$ đến $1000^\circ C$}
\loigiai{Ta có: 
	$\begin{cases}
		T_{max} = 1273\ K &\Rightarrow t_{max} = 1273 - 273 = 1\ 000^\circ C\\
		T_{min} = 263\ K C &\Rightarrow t_{min} = 263 - 273 = -10^\circ C		
	\end{cases}$.
}
\end{ex}

